\chapter{Glutamic Acid Decarboxylase (GAD) Antibodies}

\section*{What Is This Test?}
GAD antibodies are autoantibodies directed against glutamic acid decarboxylase, an enzyme found in pancreatic beta cells and the nervous system. They are markers of autoimmune beta-cell disease and can also occur in selected neurologic autoimmune syndromes.

\section*{Alternative Names}
Anti-GAD; GAD65 Antibody; Glutamic Acid Decarboxylase 65 Antibody.

\section*{Why This Test Is Ordered}
Testing helps evaluate diabetes when the type is uncertain, particularly adult-onset diabetes with lean body habitus, rapid loss of insulin production, or features suggesting latent autoimmune diabetes in adults. High levels may also be investigated in stiff-person spectrum disorders and other rare neurologic syndromes.

\section*{How This Test Is Performed}
A blood sample is tested by an immunoassay. When diabetes classification remains uncertain, additional antibodies such as IA-2, ZnT8, or insulin autoantibodies and C-peptide may be measured.

\section*{Preparation, Risks, and Limitations}
No fasting is required for the antibody test. Venipuncture is the usual risk. A negative result does not exclude autoimmune diabetes, and a low positive result may require confirmation because assay specificity and background positivity vary.

\section*{Understanding Results}
Positive GAD antibodies support autoimmune diabetes in the correct clinical setting and may predict declining endogenous insulin production. Results do not measure current glucose control and should not replace glucose, HbA1c, or C-peptide testing.

\section*{References}
American Diabetes Association standards of care on classification and diagnosis of diabetes.

\clearpage
\section*{Understanding Results and Follow-up}
The laboratory report should identify the specimen, method, units, reference interval or decision limit, and whether the result is positive, negative, abnormal, or inconclusive. Reference intervals vary with age, sex, pregnancy, specimen type, assay, and laboratory; a result outside a range is not automatically a diagnosis, and a result within a range does not always exclude disease. Discuss unexpected results with the ordering clinician, who can decide whether repeat or confirmatory testing, treatment, or referral is appropriate. Do not change medicines or delay urgent care based on an isolated result.


\section*{Regulatory Status and Authorization Date}
Regulatory status applies to the specific assay, instrument, manufacturer, and intended use rather than to the test name alone. No single approval date can be assigned to this test category. For a commercial assay, verify the current product in the relevant regulator's database: the U.S. Food and Drug Administration (FDA) clearance, approval, or authorization; India's Central Drugs Standard Control Organisation (CDSCO) licence or permission; the European Union In Vitro Diagnostic Regulation (EU IVDR) CE certificate issued through the applicable conformity-assessment route; or the United Kingdom Medicines and Healthcare products Regulatory Agency (MHRA) registration and UKCA/CE status. Laboratory-developed tests may not have an individual FDA, CDSCO, EU IVDR, or MHRA approval date.
\clearpage
